\chapter{结\texorpdfstring{\quad}{}论}
涵道式无人机凭借其结构紧凑、气动效率高、安全性好以及垂直起降能力，在低空物流、农业植保、军事侦察及应急救援等领域展现出广阔的应用前景。然而，其复杂的气动特性、近地面飞行时显著的地面效应以及强非线性的动力学耦合，给高精度飞行控制带来了严峻挑战。本文围绕涵道式无人机的气动参数在线辨识、地面效应建模与抗扰控制等关键问题，开展了系统的理论分析与实验验证，主要研究内容与结论如下：

（1）建立了涵道式无人机的完整非线性动力学模型。从机体结构出发，定义了地面坐标系与机体坐标系，采用牛顿-欧拉法推导了运动学与动力学方程，并基于现有理论分析了涵道风扇、反扭矩襟翼及控制舵面的气动力与力矩，给出了各分量的解析表达式，为后续参数辨识与控制器设计奠定了模型基础。

（2）提出了基于扩展卡尔曼滤波器的气动参数在线辨识方法。针对涵道式无人机气动参数难以离线精确标定且个体差异显著的问题，对非线性模型进行了合理简化与集总参数合并，构建了十八维增广状态空间模型，并从理论上证明了系统的局部弱可观性。数值仿真与实际飞行试验结果表明，所设计的EKF滤波器能够在典型飞行工况下实时、稳定地估计关键气动参数（如 \(K_{T0}\)、\(K_{N}\)、\(K_{cv}\) 等），估计不确定性区间有效覆盖真值，为后续非线性控制器的精确化设计提供了可靠的模型支撑。

（3）开展了地面效应下的执行器效益离线辨识研究。搭建了基于六维力/力矩传感器的变高度静态测试台架，系统测量了不同离地高度下涵道风扇推力与控制舵面偏转力矩的变化规律。采用指数衰减模型描述推力与舵面效率随无量纲高度 \(h/R\) 的变化，并通过高斯-牛顿迭代法完成了模型参数的非线性最小二乘拟合。辨识结果表明，所建模型能够准确刻画地面效应下推力的非单调变化及舵面效率的衰减特性，为近地面抗扰补偿提供了定量依据。

（4）设计了涵道式无人机抗扰控制器。首先，基于伪逆法设计了控制分配器，实现期望力/力矩向执行器指令的映射。其次，针对地面效应扰动，利用实时离地高度信息动态更新控制效率矩阵中的推力系数 \(K_{TGE}\) 与舵面偏转力系数 \(K_{cvGE}\)，实现了前馈补偿。静态台架与近地面实飞实验表明，补偿后推力与力矩的稳态误差显著降低，无人机在强地效区的高度与偏航角跟踪性能得到明显改善，稳态误差快速收敛且动态品质一致。最后，基于第三章辨识得到的气动参数，设计了非线性动态逆控制器。仿真对比结果显示，NDI控制器能够主动解耦姿态通道间的非线性耦合，消除积分项引起的动态不一致性，在姿态跟踪的快速性、稳定性和一致性方面优于串级PID控制器。

本文的研究仍存在以下不足，有待在后续工作中进一步完善：

（1）气动参数辨识实验的飞行工况覆盖范围有限。实际飞行试验主要以悬停及低速机动为主，对高速前飞、大迎角等工况的激励不足，导致部分参数（如拉力随前进比衰减系数 \(K_{TJ}\)）的收敛速度与稳态精度受到限制。后续研究应设计更丰富的飞行机动，提高参数在全飞行包线内的可辨识性。

（2）地面效应模型的适用范围存在边界。本文建立的指数衰减模型在 \(0.2 \leq h/R \leq 4.0\) 范围内拟合良好，但在极端近地（\(h/R < 0.2\)）区域，流场呈现强烈的非定常失速涡胞特性，平均推力模型的预测精度下降。后续可结合CFD仿真与动态测试，建立适用于全高度范围的混合模型。

（3）NDI控制器的实飞验证尚未完成。本文仅在仿真环境中对比了NDI与PID的控制效果，受限于样机状态与实验条件，未在实际飞行中部署NDI控制器。后续应进一步优化算法实时性，并在室内动作捕捉系统下开展NDI控制器的实飞验证。

（4）控制器的自适应能力有待提升。当前地面效应补偿依赖于离线辨识的固定参数模型，未能应对未知环境变化（如地面材质、坡度等）引起的扰动。后续可结合自适应控制或扰动观测器技术，实现在线学习与动态调整，进一步提升控制器的鲁棒性与适应性。

综上所述，本文围绕涵道式无人机的气动参数辨识、地面效应建模与抗扰控制开展了较为系统的研究，所提出的方法在理论分析与实验验证层面均取得了积极效果，为涵道式无人机在近地面复杂环境下的高精度控制提供了有效方案，并为后续更深入的研究奠定了技术基础。